\documentclass[11pt]{article}
\usepackage[T1]{fontenc}
\usepackage{amsmath,amssymb,amsthm,graphicx,hyperref}
\usepackage[margin=1in]{geometry}
\hypersetup{hidelinks}
\title{Separation-Constrained Isomap for Disconnected Manifolds\\\large Working research outline --- not a submitted manuscript}
\author{}
\date{October 2026}
\begin{document}
\maketitle
\begin{abstract}
Neighborhood graphs used in manifold learning can have multiple connected
components. Although componentwise Isomap provides a separate representation
for each component, it does not determine how these representations should
be arranged in a common space. We investigate rigid alignment subject to a
controllable geometric separation requirement, explicitly measuring the
tradeoff between preserving intracomponent geometry and preserving
intercomponent ambient distances. This document states a \emph{proposed}
framework and evaluation plan. No new theorem of substance, novelty, or
comparative empirical advantage is claimed at this stage.
\end{abstract}
\section{Research question}
Let $X=\{x_1,\ldots,x_n\}\subset\mathbb R^D$.
We study how radius-neighborhood graph connectivity affects graph-geodesic approximation
and low-dimensional representations. The radius-neighborhood rule itself is classical
and is not a proposed novelty; see \cite{tenenbaum2000}.

\section{Proposed componentwise alignment}
For each graph component $C_k$, compute graph shortest-path dissimilarities
and an Isomap representation $Z_k\subset\mathbb{R}^{d}$. Denote its row
coordinates by $z_i$. Seek a common representation
\begin{equation}
  y_i=z_i Q_k+t_k,\qquad i\in C_k,\quad
  Q_k^\top Q_k=I_d .
  \label{eq:rigid}
\end{equation}
Thus, for $i,j\in C_k$, the distance $\|y_i-y_j\|_2$ equals
$\|z_i-z_j\|_2$. This is exact preservation of the \emph{computed component
embedding}, not a guarantee that the graph accurately recovers unknown
manifold geodesics.

For $K\ge 2$, measure the fidelity of \emph{ambient} cross-component
placement through
\begin{equation}
 E_{\mathrm{cross}}(Y)=
 \left(
 \frac{
 \displaystyle\sum_{k<\ell}
 \sum_{\substack{i\in C_k\\j\in C_\ell}}
 \big(\|y_i-y_j\|_2-\|x_i-x_j\|_2\big)^2
 }{
 \displaystyle\sum_{k<\ell}
 \sum_{\substack{i\in C_k\\j\in C_\ell}}
 \|x_i-x_j\|_2^2
 }\right)^{1/2}.
 \label{eq:cross}
\end{equation}
Assume a positive denominator. Define the hull-clearance margin
$M(Y)=\min_{k<\ell}\operatorname{dist}(
\operatorname{conv}(Y_k),\operatorname{conv}(Y_\ell))$, and choose a
fixed positive reference distance $s$ so that $\mu>0$ is dimensionless.

Our proposed \emph{separation--distortion problem} is
\begin{equation}
 \begin{aligned}
 V(\mu)=\inf_{\{Q_k,t_k\}}\quad &E_{\mathrm{cross}}(Y)\\
 \text{subject to}\quad
 & y_i=z_iQ_k+t_k,\quad i\in C_k,\\
 & Q_k^\top Q_k=I_d,\quad k=1,\ldots,K,\\
 & M(Y)\ge\mu s .
 \end{aligned}
 \label{eq:frontier}
\end{equation}
For $K=1$, cross-component stress and clearance are not applicable.
Ordered projections onto a common axis with positive gaps yield a simple
\emph{sufficient} constraint for hull separation, but impose a stricter
feasible region than unrestricted Euclidean hull clearance.

\paragraph{Elementary mathematical properties.}
Rigid transformations preserve all intracomponent distances.
Increasing $\mu$ shrinks the feasible set in \eqref{eq:frontier},
so the \emph{global} optimal value $V(\mu)$ is nondecreasing (with
$+\infty$ allowed for infeasibility). This does not guarantee global
convergence or monotonicity of values returned by a local optimizer.
These simple properties are proved in the repository note
\texttt{PROP-0002-rigid-separation-frontier.md}; they are not claimed
as new mathematical results.

\section{Required evaluation and limitations}
Before any performance claim, implement an unsupervised selection of
neighborhood radius and verify how often graph components correspond to
true generating manifolds. Vary the required normalized clearance
$\mu$ and report (i) within-manifold intrinsic-distance stress where
ground truth is known, (ii) the achieved hull clearance and overlap
rate, and (iii) cross-component ambient-distance stress.
Graph-target stress should be distinguished from true geodesic error.

Use identical sampled data, independent random seeds, tuned PCA/Isomap
and relevant multi-manifold/rigid-alignment methods, fixed evaluation
pairs, and explicit policies for disconnected baseline graphs.
Compare up to six rotated manifolds, and vary sampling density,
orientation, noise, and original separation. Record convergence,
computational cost, and the effect of replacing full cross-pair
objectives by representative anchor approximations.

\paragraph{Two main figures (proposed).}
First, a controlled 3D example and 2D projections comparing component
overlap and intrinsic fidelity. Second, a separation--distortion
frontier at varying clearance, with independent-sample variation
and competitive baselines. The existing EXP-0002--EXP-0008 experiments
are local archive work and are \emph{not yet merged or reproduced} on
the repository's current main branch. No numerical claims are made
in this outline.

\paragraph{Novelty boundary.}
Proximity graphs, multi-manifold embedding, rigid alignment, and
separation metrics are existing ideas. A publication requires either
a demonstrably different and useful constrained method, a nontrivial
new analysis, or strong reproducible improvements over the relevant
prior art. Visual non-overlap alone is insufficient.

\bibliographystyle{plain}
\bibliography{../literature/references/library}
\end{document}
